\subsection{Larger Linear Systems} 
We have already studied systems of \makeblanksmall linear equations with \makeblanksmall variables. We can add more equations and more variables.
\begin{itemize}
\item If a system has more equations than variables, the system is \term{overdetermined}, and will typically have \makeblank[1in] solutions. An overdetermined system has \makeblank[1.5in] information.
\end{itemize}
\begin{Example}
Consider the system
\begin{syspic}{5}
\draw (-.5\textwidth,0) node{%
$\begin{systemL}
y=x-2&\\
y=-\frac{1}{2}x+1&\\
y=2x-1&
\end{systemL}$};
\draw[graph] (-3,-5) -- (5,3);
\draw[graph] (-2,5) -- (5,3/2);
\draw[graph] (-2,-5) -- (3,5);
\end{syspic}
Every \makeblank[1in] of equations has a shared solution, but there is no solution shared by \makeblank[1.5in].
\end{Example}
\begin{itemize}
\item If a system has fewer equations than variables, the system is \term{underdetermined}, and will typically have \makeblank[1in] solutions. An underdetermined system has \makeblank[1.5in] information.
\end{itemize}
\begin{Example}
Consider the system
\begin{syspic}{5}
\draw (-.5\textwidth,0) node{%
$\begin{systemL}
5x+2y=8&
\end{systemL}$};
\draw[graph] (-0.4,5) -- (3.6,-5);
\end{syspic}
Every point on the \makeblank[1in] is a solution to this ``system.''
\end{Example}
\begin{itemize}
\item To have a unique solution, we must typically have a \term{balanced} system, in which the number of equations is \makeblank[1in] the number of variables.
\end{itemize}
In the rest of this section, we will concern ourselves with systems of \makeblanksmall linear equations in \makeblanksmall variables.
\begin{Note}
The ``typically'' refers to the fact that well chosen systems of each type may behave differently. If we write down a system at random, it will almost certainly have the typical behavior.
\end{Note}